\scanpage{100}

d'où enfin une application

\[
\begin{aligned}
 Z&\longrightarrow Z'_\Lambda=\Lambda^*/\{\pm1_\Lambda\}\\
 \lambda&\longmapsto(\lambda_{\Gamma,s})^{\unsure{\natural|\sharp|4}}
\end{aligned}\tag*{(66)}
\]

qui a priori dépend du choix de $i_0\in\Pi$, et de $s\in S$. Mais on trouve

\[
\begin{aligned}
 &\text{L'application précédente est indépendante}\\
 &\text{du choix de }s\in S\text{ et de}\\
 &i_0\in\Pi,\ \text{et c'est un \uline{isomorphisme}}\\
 &\text{de groupes (en fait, c'est l'}\\
 &\text{isomorphisme inverse de l'isomorphisme}\\
 &\text{obtenu p. }47'\text{).}
\end{aligned}\tag*{(67)}
\]

% slanted
\margin{à vérifier}

Ainsi, j'ai explicité, en termes d'une quelconque des structures $P_{M,i}$ -- ou de l'ens. de ces structures, au choix \inserted{ou de la structure $F_M$ \unsure{seulement}} -- l'opération de $Z'_\Lambda$ sur $S$, et par là, l'opération de $Z'_\Lambda$ sur $\Pi$, qui fait de $\Pi$ un $Z'_\Lambda$-torseur.

Explicitons maintenant le $\Lambda^*$-torseur $\vec\Pi^{*}$, on a une bijection canonique

\[
\vec\Pi\simeq\coprod_{i\in\Pi}\bigl[\underbrace{\operatorname{Or}(P_{\deleted{\unsure{M}},i})}_{\deleted{\text{\unsure{ou plutôt ens. des deux orientations de |petit ens. des deux orientations de }P_i}}}=\underbrace{\mu(i)}_{\text{ens des deux or de $P_i$}}\bigr]\tag*{(68)}
\]
% archive sheet number: 99
